Synthetic data generated by large language models is increasingly used for fine-tuning, yet it suffers from reduced sample efficiency and, when applied recursively, causes model collapse. Despite growing empirical evidence of these failures, the underlying optimization-level mechanism remains poorly understood. We hypothesize that because LLM-generated text closely matches the model's own distribution, it produces both smaller and directionally narrower gradients during training, starving the model of learning signal for distributional tails. We test this hypothesis through four controlled experiments using GPT-2 Small fine-tuned on WikiText-2. We find that synthetic data has $3\times$ lower median perplexity than human data, produces 15\% smaller gradient norms ($p < 10^{-45}$), and---most strikingly---concentrates gradient signal into just 1--2 principal components compared to 51+ for human data. This directional collapse means the model reinforces already-dominant patterns while receiving virtually no gradient signal for rare or diverse linguistic phenomena. The practical consequence is a $1.4$--$1.8\times$ reduction in sample efficiency. Our results provide the first empirical gradient-level mechanism linking low perplexity to model collapse: it is not merely that synthetic gradients are smaller, but that they point in far fewer directions.
